\section{The original three special points and an exact real torus quotient}
\label{sec:s6q-three-points}

The input for this section is the three-point period and gluing data in
the supplied unpublished manuscript \emph{A compact complex threefold
fibred by tori over the projective line, and the six-sphere}, its independent
TeX reconstruction and its accompanying workbench proofs. The precise
source loci are the reconstruction's Section~3 and
\texttt{candidate\_geometry.tex},
\texttt{exact\_algebra.tex}, and \texttt{integral\_leray.tex} in that
workbench. We use their explicitly marked local period family and overlap
formulas as source input. No identification of the global source manifold
with a sphere is used in the following quotient calculation.

\subsection{The retained base points, periods and real marking}

The original base and its distinguished points are
\[
 B=\mathbb P^1_t,\qquad
 p_1=0,\quad p_2=1,\quad p_0=\infty,\qquad
 t_c=1/t,\qquad B^\circ=B\setminus\{p_0,p_1,p_2\}.
\]
The finite orders are $m_1=3$ and $m_2=4$. On the marked orbifold cover
$\pi:\mathfrak h_z\to B\setminus\{p_0\}$, put
$g_0=(g_1g_2)^{-1}$. The chosen lifts $z_j$ of $p_j$ have coordinates
\[
 s_j(z)=\frac{z-z_j}{z-\overline{z_j}},\qquad
 s_j(g_jz)=e^{-2\pi i/m_j}s_j(z),\qquad
 t_j\circ\pi=s_j^{m_j}.
\]
The inverse of the displayed Cayley coordinate is
$z=(z_j-s_j\overline{z_j})/(1-s_j)$. The modular period is the separate map
$\tau:\mathfrak h_z\to\mathfrak h_\tau$, with
\[
 j(\tau)=1728(t\circ\pi),\qquad
 \tau(z_1)=\rho=e^{\pi i/3},\qquad \tau(z_2)=i,
\]
\[
 \mu(z_1)=\frac{2-\rho}{3},\qquad \mu(z_2)=\frac{1-i}{2}.
\]
Its exact local forms are
$\tau-\rho=s_1U_1(s_1)$ and $\tau-i=s_2^2U_2(s_2)$ with
$U_1(0)U_2(0)\ne0$. These units and the local base-coordinate changes
$t=t_1V_1(t_1)$ and $t-1=t_2V_2(t_2)$ remain part of the original data;
$V_j(0)\ne0$ does not assign them the value one.

The ordered source lattices and periods are
\[
 V=\mathbb Z\langle\gamma,u,w,\delta\rangle,\qquad
 \Lambda=V^*=
 \mathbb Z\langle\widehat\gamma,\widehat u,\widehat w,\widehat\delta\rangle,
\]
\[
 \Pi(z)=[Z(z)\mid I_2],\qquad
 Z(z)=\begin{pmatrix}6\mu(z)&\tau(z)\\\beta(z)&\mu(z)\end{pmatrix}.
\]
Vectors are columns in this ordered basis. The real map underlying the
period map is exactly
\begin{equation}
 \Pi(z)(a,b,c,d)
   =(6\mu a+\tau b+c,\ \beta a+\mu b+d).
 \label{s6q:period-map}
\end{equation}
The source admissible period parameter ensures
\[
 \operatorname{Im}\tau>0,\qquad
 D=\operatorname{Im}\beta-
       \frac{6(\operatorname{Im}\mu)^2}{\operatorname{Im}\tau}<0.
\]
Thus the real determinant in the output order
$(\operatorname{Re}\zeta_1,\operatorname{Im}\zeta_1,
  \operatorname{Re}\zeta_2,\operatorname{Im}\zeta_2)$ is
$(\operatorname{Im}\tau)D\ne0$. The admissible constant in $\beta$ is
retained; no special numerical value is selected here. For an ordinary
lifted base open set $U$, the marking is the real analytic isomorphism
\begin{equation}
 \mathcal M_U:U\times(\mathbb R^4/\mathbb Z^4)
       \longrightarrow \mathcal T|_U,
 \qquad (z,x+\mathbb Z^4)\longmapsto[z,\Pi(z)x].
 \label{s6q:marking}
\end{equation}
The inverse uses the real inverse of \eqref{s6q:period-map}. This keeps
the full source complex torus in the domain of the next map.

The original inverse-transpose monodromies on $\Lambda$ are
\begin{equation}
 \begin{split}
 A_1&=\begin{pmatrix}
 1&0&0&0\\6&0&1&0\\-6&-1&-1&0\\-2&1&0&1
 \end{pmatrix},\qquad
 A_2=\begin{pmatrix}
 1&0&0&0\\0&0&-1&0\\-6&1&0&0\\3&0&1&1
 \end{pmatrix},\\
 M_0&=\begin{pmatrix}
 1&0&0&0\\0&1&0&0\\0&1&1&0\\-1&0&0&1
 \end{pmatrix}.
 \end{split}
 \label{s6q:full-matrices}
\end{equation}
The source relations are $A_1^3=A_2^4=I_4$ and $A_1A_2M_0=I_4$.
The two original fixed twist vectors are
\begin{equation}
 v_1=(1,2,-4,0)^{\mathsf T},\qquad
 v_2=(-1,-3,3,0)^{\mathsf T},\qquad A_jv_j=v_j.
 \label{s6q:full-twists}
\end{equation}
In particular the fourth marked coordinate is present in every original
matrix and twist before taking a quotient.

\subsection{The exact circle quotient and its section defects}

Let $e_1,e_2,e_3,e_4$ denote the four vectors of the original ordered
real marking, and use $\bar e_1,\bar e_2,\bar e_3$ in the quotient.
Define
\[
 P:\mathbb R^4\longrightarrow\mathbb R^3,
 \quad P(a,b,c,d)=(a,b,c),\qquad
 \iota:\mathbb R^3\longrightarrow\mathbb R^4,
 \quad\iota(a,b,c)=(a,b,c,0).
\]
Their matrices are $P=[I_3\mid0]$ and
$\iota=[I_3\mid0]^{\mathsf T}$. Since $P\mathbb Z^4=\mathbb Z^3$,
they induce smooth homomorphisms
\begin{equation}
 \overline P:\mathbb T^4\longrightarrow\mathbb T^3,
 \quad[x]\longmapsto[Px],\qquad
 \overline\iota:\mathbb T^3\longrightarrow\mathbb T^4,
 \quad[y]\longmapsto[\iota y],
 \qquad \mathbb T^n=\mathbb R^n/\mathbb Z^n.
 \label{s6q:torus-quotient}
\end{equation}

\begin{proposition}[Typed quotient, kernel and full monodromy intertwining]
The maps in \eqref{s6q:torus-quotient} give the split exact sequence of
marked compact real Lie groups
\begin{equation}
 0\longrightarrow\mathbb Re_4/\mathbb Ze_4
   \longrightarrow\mathbb T^4
   \xrightarrow{\ \overline P\ }\mathbb T^3\longrightarrow0,
 \qquad \overline P\,\overline\iota=\operatorname{id}_{\mathbb T^3}.
 \label{s6q:exact-sequence}
\end{equation}
The quotient monodromies in the original first three coordinates are
\begin{equation}
 \bar A_1=\begin{pmatrix}1&0&0\\6&0&1\\-6&-1&-1\end{pmatrix},\qquad
 \bar A_2=\begin{pmatrix}1&0&0\\0&0&-1\\-6&1&0\end{pmatrix},\qquad
 \bar M_0=\begin{pmatrix}1&0&0\\0&1&0\\0&1&1\end{pmatrix}.
 \label{s6q:quotient-matrices}
\end{equation}
They satisfy
\[
 PA_j=\bar A_jP,\quad PM_0=\bar M_0P,\quad
 \bar A_1^3=\bar A_2^4=I_3,\quad \bar A_1\bar A_2\bar M_0=I_3.
\]
The section has the exact defects
\begin{equation}
 \begin{aligned}
 (A_1\iota-\iota\bar A_1)(a,b,c)&=(-2a+b)e_4,\\
 (A_2\iota-\iota\bar A_2)(a,b,c)&=(3a+c)e_4,\\
 (M_0\iota-\iota\bar M_0)(a,b,c)&=-ae_4.
 \end{aligned}
 \label{s6q:section-defects}
\end{equation}
\end{proposition}

\begin{proof}
For $[x]\in\ker\overline P$, the first three entries of $x$ are
integers. Subtracting that vector of $\mathbb Z^4$ leaves precisely
$de_4$. Two such representatives give the same class exactly when their
fourth coordinates differ by an integer. This proves the stated kernel.
Surjectivity and the section identity follow by applying $P$ to
$\iota(a,b,c)$. The same equations show that
$\mathbb T^4/(\mathbb Re_4/\mathbb Ze_4)\to\mathbb T^3$,
$[[x]]\mapsto[Px]$, is bijective with inverse $[y]\mapsto[[\iota y]]$.
Both maps are smooth: in the ordinary local covering coordinates they are
the displayed linear projection and inclusion followed by the quotient.

Each matrix in \eqref{s6q:full-matrices} fixes $e_4$, and its first three
rows have zero fourth entry. Taking those rows proves all three
intertwiners and the matrices in \eqref{s6q:quotient-matrices} by literal
multiplication. Applying $P$ to the source product and power identities,
and then composing with the surjection $P$, proves the corresponding
identities in dimension three. The lower row of each full matrix applied
to $(a,b,c,0)$ gives exactly \eqref{s6q:section-defects}. These defects
are generally nonzero even as circle classes; thus the section is a
right inverse of the marked quotient without being a monodromy-equivariant
section.
Indeed there is no real linear section equivariant even for $M_0$ alone.
Any linear section has the form $\iota_L(y)=(y,Ly)$ for a row $L$.
Equivariance would require
$L(\bar M_0-I_3)=(-1,0,0)$. Its left side is $(0,L_3,0)$,
contradicting the first coordinate of the required equality.
\end{proof}

On the source fibre, this gives the concrete map
\begin{equation}
 p_z:\mathbb C^2/\Pi(z)\Lambda\longrightarrow\mathbb T^3,
 \qquad [\zeta]\longmapsto[P\Pi(z)^{-1}\zeta],
 \label{s6q:actual-fibre-map}
\end{equation}
where the inverse is real linear. Its kernel is the original circle
$\mathbb R\Pi(z)e_4/\mathbb Z\Pi(z)e_4$; here
$\Pi(z)e_4=(0,1)^{\mathsf T}$ exactly. The map
$[y]\mapsto[\Pi(z)\iota y]$ is its marked right inverse.
Equation \eqref{s6q:marking} proves the smooth dependence of the quotient
and section on the original base variable. This construction concerns a
three-real-dimensional fibre quotient. Keeping the complex base produces
a five-real-dimensional family; no identification of that total space
with a three-dimensional physical domain is implicit.

\subsection{The finite affine actions and the complete logarithmic overlaps}

Put
\begin{equation}
 c_1=P(v_1/3)=\tfrac13(1,2,-4)^{\mathsf T},\qquad
 c_2=P(v_2/4)=\tfrac14(-1,-3,3)^{\mathsf T}.
 \label{s6q:finite-directions}
\end{equation}
For $j=1,2$ the full marked map
$L_j(x)=A_jx+v_j/m_j$ on $\mathbb T^4$ descends to
\[
 \bar L_j([y])=[\bar A_jy+c_j],\qquad
 \overline P L_j=\bar L_j\overline P.
\]
All matrices and translations are those obtained above. Since
$\bar A_jc_j=c_j$, induction gives
\begin{equation}
 \bar L_j^{k}([y])=[\bar A_j^{k}y+kc_j].
 \label{s6q:affine-powers}
\end{equation}
At $k=m_j$ the matrix is the identity and the translation is the integer
vector $Pv_j$, so the map is the identity on $\mathbb T^3$.
The first coordinate of every $\bar A_j^{k}y$ is $y_1$.
For $0<k<3$, the first-coordinate translation of $\bar L_1^{k}$ is
$k/3\notin\mathbb Z$; for $0<k<4$, that of $\bar L_2^{k}$ is
$-k/4\notin\mathbb Z$. Thus every nonidentity power is fixed-point free,
and the two affine torus actions have exactly orders three and four.
On $\Delta_j\times\mathbb T^3$ the corresponding maps are
$(z,[y])\mapsto(g_jz,\bar L_j[y])$, and the same argument proves freeness
also over the fixed base point. The quotient fibre over $p_j$ is therefore
the smooth compact real quotient $\mathbb T^3/\langle\bar L_j\rangle$.

The overlap is a varying translation whose complete real expression must
also be retained. Define
\[
 J_z=\Pi(z)^{-1}i\Pi(z):\mathbb R^4\to\mathbb R^4,
 \qquad \alpha(z)=\frac{\log s_j(z)}{2\pi i},
\]
\begin{equation}
 a_j(z)=\operatorname{Re}\alpha(z)\,v_j+
          \operatorname{Im}\alpha(z)\,J_zv_j,
 \qquad \bar a_j(z)=Pa_j(z).
 \label{s6q:full-overlap-translation}
\end{equation}
Here is an explicit coordinate formula for the full $J_z$, retaining all
entries of the period matrix. Write $\mu_R=\operatorname{Re}\mu$,
$\mu_I=\operatorname{Im}\mu$, and likewise for $\tau,\beta$. For
$x=(a,b,c,d)$ set
\[
 p=6\mu a+\tau b+c,\qquad q=\beta a+\mu b+d,
 \qquad E=6\mu_I^2-\tau_I\beta_I=-\tau_ID\ne0.
\]
Then $J_zx=(a',b',c',d')$ is exactly
\begin{equation}
 \begin{aligned}
 a'&=\frac{\mu_I\operatorname{Re}p-\tau_I\operatorname{Re}q}{E},&
 b'&=\frac{-\beta_I\operatorname{Re}p+6\mu_I\operatorname{Re}q}{E},\\
 c'&=-\operatorname{Im}p-6\mu_Ra'-\tau_Rb',&
 d'&=-\operatorname{Im}q-\beta_Ra'-\mu_Rb'.
 \end{aligned}
 \label{s6q:complex-structure-coordinates}
\end{equation}
Indeed the imaginary parts of $\Pi(z)J_zx=i(p,q)$ are the two linear
equations for $a',b'$ with determinant $E$; its real parts give $c',d'$.
For the two actual twist inputs one uses, without changing their signs,
\[
 (p,q)_{v_1}=(6\mu+2\tau-4,\beta+2\mu),\qquad
 (p,q)_{v_2}=(-6\mu-3\tau+3,-\beta-3\mu).
\]
Equations \eqref{s6q:full-overlap-translation} and
\eqref{s6q:complex-structure-coordinates} consequently specify every
coordinate of the descended translation.

The source period covariance gives
$J_{g_jz}A_j=A_jJ_z$. Clockwise continuation of the chosen logarithm
changes $\alpha$ by $-1/m_j$, so
\begin{equation}
 a_j(g_jz)=A_ja_j(z)-v_j/m_j,
 \qquad \bar a_j(g_jz)=\bar A_j\bar a_j(z)-c_j.
 \label{s6q:overlap-covariance}
\end{equation}
To verify the first identity directly, the real part of $\alpha$ acquires
the displayed subtraction; its imaginary part stays fixed, and
$A_jv_j=v_j$ gives $J_{g_jz}v_j=A_jJ_zv_j$. Applying the proved
intertwiner $P$ gives the second identity.

On the punctured lifted family put
\[
 \bar h_j(z,[y])=(z,[y+\bar a_j(z)]).
\]
Changing $\log s_j$ by $2\pi in$ changes $\bar a_j$ by the integer
vector $nPv_j$, proving branch independence. Its inverse is translation
by $-\bar a_j(z)$. Applying \eqref{s6q:overlap-covariance} to the three
successive maps gives the exact conjugacy
\begin{equation}
 \bar h_j\,(g_j,\bar L_j)\,\bar h_j^{-1}
    :(z,[y])\longmapsto(g_jz,[\bar A_jy]).
 \label{s6q:overlap-conjugacy}
\end{equation}
The cancelled translation is
$-\bar A_j\bar a_j(z)+c_j+\bar a_j(g_jz)=0$.
This proves descent of the complete logarithmic gluing to the real
circle quotient, including its inverse and every period-dependent
coefficient. Circle cosets are preserved because the full linear maps
fix $e_4$ and fibre translations preserve cosets. No preservation of a
chosen Euclidean or Hermitian metric is required or asserted.

There is a concrete reason to retain these overlap translations when
combining the three source points. In a single uncorrected marked fibre,
compose the two finite affine formulas and the linear cusp formula. Their
linear product is $\bar A_1\bar A_2\bar M_0=I_3$, but their translation is
\[
 c_1+\bar A_1c_2=(1/12,-1/12,1/6)^{\mathsf T}\notin\mathbb Z^3.
\]
Thus those three uncorrected affine formulas are not the three meridians
in the same middle-family marking. The exact marking change is
\eqref{s6q:overlap-conjugacy}; it identifies each finite formula with its
linear middle-family representative, where the product relation holds.
This computes the discrepancy and the maps correcting it.

\subsection{The original cusp direction and three independent shear data}

The cusp's complete period expression is
\[
 e^{2\pi i\tau}=t_cu_1(t_c),\quad u_1(0)=1/1728,\quad
 h=(2\pi i)^{-1}\log u_1,\quad s=\tau-h,\quad b=\beta+\tau,
\]
\begin{equation}
 Z=sB_0+C(t_c),\qquad B_0=\begin{pmatrix}0&1\\-1&0\end{pmatrix},\qquad
 C(t_c)=\begin{pmatrix}6\mu(t_c)&h(t_c)\\b(t_c)-h(t_c)&\mu(t_c)\end{pmatrix}.
 \label{s6q:cusp-periods}
\end{equation}
The full holomorphic unit $u_1$ and all entries of $C$ remain present.
Define the original cusp sublattice and quotient with their ordered bases:
\[
 \Lambda_{\rm tor}=\mathbb Z\langle e_3,e_4\rangle,\qquad
 \bar\Lambda=\Lambda/\Lambda_{\rm tor}
   =\mathbb Z\langle[e_1],[e_2]\rangle.
\]
On the original marked lattice,
\[
 (M_0-I)(a,b,c,d)=(0,0,b,-a),\qquad
 \ker(M_0-I)=\operatorname{im}(M_0-I)
   =\mathbb Z\langle e_3,e_4\rangle.
\]
The quotient sends this cusp image lattice onto $\mathbb Z\bar e_3$,
with kernel $\mathbb Ze_4$. In particular the unchanged source vector
$e_2$ supplies
\begin{equation}
 c_0=P((M_0-I)e_2)=\bar e_3=(0,0,1)^{\mathsf T}.
 \label{s6q:cusp-direction}
\end{equation}
The other primitive image $(M_0-I)e_1=-e_4$ is killed by this precise
quotient. This records the lost cusp direction explicitly. Translation
by the integer vector $c_0$ is the identity on $\mathbb T^3$; the same
vector defines a nonzero tangent direction there. These two statements
follow respectively from taking its class modulo $\mathbb Z^3$ and from
the differential of the covering $\mathbb R^3\to\mathbb T^3$.

The relation with the original two-dimensional cusp torus is therefore
the following exact restriction. Denote its full marked torus by
\[
 \mathbb T^2_{\rm tor}=
 (\mathbb Re_3\oplus\mathbb Re_4)/(\mathbb Ze_3\oplus\mathbb Ze_4),
 \qquad K_4=\mathbb Re_4/\mathbb Ze_4,\quad
 K_3=\mathbb R\bar e_3/\mathbb Z\bar e_3.
\]
In addition to \eqref{s6q:exact-sequence}, one has
\begin{equation}
 \begin{gathered}
 0\longrightarrow K_4\longrightarrow\mathbb T^2_{\rm tor}
 \xrightarrow{\ [re_3+ue_4]\mapsto[r\bar e_3]\ }
 K_3\longrightarrow0,\\
 0\longrightarrow\mathbb Re_4\longrightarrow
 \mathbb Re_3\oplus\mathbb Re_4
 \xrightarrow{\ (r,u)\mapsto r\ }
 \mathbb R\bar e_3\longrightarrow0.
 \end{gathered}
 \label{s6q:cusp-torus-restriction}
\end{equation}
The vertical inclusions of the middle and final terms into $\mathbb T^4$
and $\mathbb T^3$ make a commuting diagram with
\eqref{s6q:exact-sequence}: on every representative,
$P(re_3+ue_4)=r\bar e_3$. The same equality proves commutation of the
differential sequences. Their kernels, surjections and sections
$[r]\mapsto[re_3]$ follow by testing the displayed coordinates and their
integer lattices. The original $B_0$ is also retained in this diagram:
the map $\bar\Lambda\to\Lambda_{\rm tor}$ sends
$(a,b)\mapsto be_3-ae_4$, and composition with $P$ is
$(a,b)\mapsto b\bar e_3$. This proves the exact relationship between
the two real cusp coordinates, their two-torus, and the one-dimensional
image inside the three-dimensional marked quotient.

The cusp punctured-overlap map on the original complex covering is
\[
 (z,\zeta_1,\zeta_2)\longmapsto
 (e^{2\pi i\zeta_1},e^{2\pi i\zeta_2},e^{2\pi is(z)}).
\]
In these coordinates the $e_4$ circle acts by
$(x_1,x_2,t_c)\mapsto(x_1,e^{2\pi i r}x_2,t_c)$,
$r\in\mathbb R/\mathbb Z$. This follows by exponentiating
$\Pi(z)e_4=(0,1)$. It commutes with the full toric multipliers
\[
 (x,t_c)\longmapsto
 \bigl(e^{2\pi iC(t_c)\bar\lambda}
                t_c^{B_0\bar\lambda}x,t_c\bigr),
 \qquad\bar\lambda\in\mathbb Z^2.
\]
Thus the overlap identifies the same circle quotient on the punctured
cusp model. The calculation makes no assertion that this circle action
remains free on the singular central fibre; the defined three-dimensional
marked fibre quotient is on the smooth torus locus.

Finally, in the original order $(p_1,p_2,p_0)$ the three directions form
\begin{equation}
 [c_1\ c_2\ c_0]
 =\begin{pmatrix}
 1/3&-1/4&0\\2/3&-3/4&0\\-4/3&3/4&1
 \end{pmatrix},\qquad
 \det[c_1\ c_2\ c_0]=-\frac1{12}.
 \label{s6q:direction-determinant}
\end{equation}
The integer wave covectors in the same order are
\begin{equation}
 k_1=(4,0,1),\qquad k_2=(3,0,1),\qquad k_0=(0,1,0),
 \qquad
 \det\begin{pmatrix}k_1\\k_2\\k_0\end{pmatrix}=-1.
 \label{s6q:wave-covectors}
\end{equation}
Direct multiplication gives $k_jc_j=0$ for all three labels: explicitly
$(4-4)/3=0$, $(-3+3)/4=0$, and $0=0$. The first determinant is the
upper-left minor $-1/4+1/6=-1/12$; expansion along the last row of the
second matrix gives $-(4-3)=-1$. Hence the source directions span the
whole quotient tangent space and the wave covectors form an integral
basis of its dual lattice. The direction map
$(r_1,r_2,r_0)\mapsto r_1c_1+r_2c_2+r_0c_0$ has the exact inverse matrix
\[
 [c_1\ c_2\ c_0]^{-1}
   =\begin{pmatrix}9&-3&0\\8&-4&0\\6&-1&1\end{pmatrix}.
\]
Multiplying this displayed matrix on either side by
\eqref{s6q:direction-determinant} gives $I_3$, proving the inverse in
the original coordinates. For each $j$, each scalar amplitude $a_j$ and
each smooth one-periodic function $F_j$, the literal quotient field
$y\mapsto a_jc_jF_j(k_jy)$ is periodic and has divergence
$a_j(k_jc_j)F_j'(k_jy)=0$. This final derivative identity uses the
explicit quotient coordinates; constructing and checking its subsequent
Navier--Stokes dynamics is a further calculation on these defined fields.
These individual modes are fields on the marked $\mathbb T^3$. For
clarity, their first two wave covectors obey
$k_1\bar A_1=(-2,-1,-1)$ and $k_2\bar A_2=(-3,1,0)$, whereas
$k_0\bar M_0=k_0$. Therefore no automatic descent of the individual
finite-point modes to $\mathbb T^3/\langle\bar L_j\rangle$ is asserted.
There is an explicit field map to that quotient: for a smooth
one-periodic real function $F$, the finite orbit sum
\[
 \mathcal R_jF(y)=c_j\sum_{\ell=0}^{m_j-1}
                         F(k_j\bar A_j^\ell y)
\]
is invariant as a vector field under $\bar L_j$. Indeed
$\bar A_jc_j=c_j$ and
$k_j\bar A_j^\ell c_j=k_jc_j=0$; substitution of
$\bar A_jy+c_j$ cycles the sum and gives
$\mathcal R_jF(\bar L_jy)=\bar A_j\mathcal R_jF(y)$.
Each summand has zero divergence by the same coordinate derivative
calculation. The local quotient covering then gives a unique smooth
field downstairs: at $[y]$ use its differential on $\mathcal R_jF(y)$;
the displayed covariance proves independence of the lift. This supplies
the exact field descent map while keeping its differential-operator
and metric questions separate for the subsequent dynamics calculation.
